% gerado por scripts/build_paper_assets.py -- nao editar a mao
\begin{table}[htbp]
\centering
\small
\setlength{\tabcolsep}{4pt}
\caption{Expanding versus sliding 60-month training window, evaluated on identical
test origins. A negative difference means the expanding window is more accurate.}
\label{tab:janela}
\begin{tabular}{lrrrcrc}
\toprule
Model & Expanding & Sliding 60 & Diff. (pp) & 95\% CI & $p$ & DM cells $p<0.05$ \\
\midrule
Prophet & 4.70 & 5.39 & $-0.69$ & $[-1.159, -0.219]$ & 0.005 & 0 of 6 \\
SARIMA & 4.80 & 5.46 & $-0.66$ & $[-1.008, -0.346]$ & 0.000 & 3 of 6 \\
TimesFM & 4.83 & 4.74 & $+0.09$ & $[-0.125, +0.307]$ & 0.424 & 0 of 6 \\
CatBoost & 6.59 & 6.64 & $-0.05$ & $[-0.288, +0.176]$ & 0.661 & 0 of 6 \\
XGBoost & 6.88 & 6.86 & $+0.02$ & $[-0.258, +0.303]$ & 0.879 & 0 of 6 \\
\midrule
CatBoost, enriched & 6.09 & 6.59 & $-0.50$ & --- & --- & --- \\
XGBoost, enriched & 6.39 & 6.79 & $-0.41$ & --- & --- & --- \\
\bottomrule
\end{tabular}

\vspace{0.5em}
\begin{minipage}{\textwidth}
\footnotesize
The only difference between the two runs is how much past each model may use: the
expanding window trains from the start of the series to the origin (60 to 162 months),
the sliding window on the most recent 60 months only. Test origins, horizons and
evaluation are identical, so the bootstrap is paired by window. Under the pre-declared
criterion, only SARIMA meets both conditions. The two rows below the second rule are the best feature engineering variant of each boosting model (seasonal differencing, calendar terms, one model per horizon; Table~\ref{tab:variantes}), run under both window policies. They lose 0.41 to 0.50 percentage points under the short window, so the advantage of feature engineering does not survive a short history either. Their interval, $p$ and DM cells are left blank because only the aggregate metrics of the sliding run were retained, not its individual forecasts, and the paired bootstrap needs the forecasts.
\end{minipage}
\end{table}
